Ce manuscrit \'etudie une approche texte d'abord pour la d\'etection de moments saillants dans de longues vid\'eos. Au lieu de confier la d\'ecision finale \`a une inf\'erence multimodale directe sur la vid\'eo, nous transformons le contenu en un \'etat textuel structur\'e et persistant, stock\'e en base de donn\'ees, puis nous comparons plusieurs vues textuelles sur le jeu d'\'evaluation \texttt{validation29}. Les r\'esultats montrent que des repr\'esentations textuelles surpassent l'entr\'ee vid\'eo directe dans la comparaison contr\^ol\'ee sous \texttt{Gemini 3.1 Flash Lite}, tout en co\^utant moins cher \`a l'inf\'erence finale. Ils montrent aussi que la textualisation visuelle actuelle n'aide pas automatiquement et que le choix entre enrichissement de contexte et ajustement fin d\'epend de la taille du mod\`ele. La contribution centrale est donc double: une comparaison reproductible entre repr\'esentations texte et vid\'eo directe, et une repr\'esentation vid\'eo transform\'ee en texte qui reste lisible, r\'eutilisable et exploitable dans le temps.
